\begin{table}[htbp]
\centering
\caption{Diagnostics for the baseline value-relevance model (Model 1)}
\label{tab:t6_diagnostics}
\begin{tabular}{lrrr}
\toprule
Diagnostic & Statistic & p / threshold & Note \\
\midrule
Model 1 VIF (EPS) & 1.16 & - & tolerance 0.862 \\
Model 1 VIF (BVPS) & 1.16 & - & tolerance 0.862 \\
Model 2 VIF (EPS) & 10.33 & - & tolerance 0.097 \\
Model 2 VIF (BVPS) & 2.28 & - & tolerance 0.439 \\
Model 2 VIF (EPS\_POST) & 15.15 & - & tolerance 0.066 \\
Model 2 VIF (BVPS\_POST) & 24.13 & - & tolerance 0.041 \\
Breusch-Pagan LM & 14.37 & p = 0.213 & Koenker robust form; statsmodels robust=True \\
Classical BP LM & 28.18 & p = 0.003 & Normality-based classical variant \\
Max Cook's D & 0.552 & 4/n = 0.143 & 5 of 28 observations exceed the threshold \\
\midrule
\multicolumn{4}{l}{\footnotesize Notes: VIFs are computed on the substantive regressors only (fixed-effect dummies excluded); Model 2 includes both interaction terms. The first LM row is the Koenker robust form and the second is classical Breusch-Pagan. HC3 changes the covariance estimate; it does not remove influential observations. See the Q-Q and leave-one-out diagnostics in F5 and Appendix E.} \\
\bottomrule
\end{tabular}
\end{table}
